\subsection{Operator Functions}


% **Role:** structure revealed by invariant subspaces and eigenbehavior.
%
% Finite spectral calculus, positive operators, and spectral commutativity.

This subsection develops finite spectral functions of normal operators and positive square roots.\\
It then characterizes commutativity through spectral projections and a single normal generator.

\begin{definition} \label{def:normal-spectral-resolution-notation} \hlt{Spectral Resolution Notation}\\
Let $\mathcal{V}$ be a finite-dimensional complex inner-product space, and let $T\in\LL(\mathcal{V})$ be normal. The associated spectral notation is:
\begin{enumerate}[label=(\roman*)]
\setlength{\itemsep}{0pt}
\item \hlt{Spectral Set}: $\Omega_T\defeq\operatorname{spec}(T)$.
\item \hlt{Spectral Projections}: For $\lambda\in\Omega_T$, $P_{\lambda}$ is the orthogonal projection onto $E_{\lambda}(T)$.
\item \hlt{Spectral Resolution}: $(P_{\lambda})_{\lambda\in\Omega_T}$, with $T=\sum\limits_{\lambda\in\Omega_T}\lambda P_{\lambda}$ and $I_{\mathcal{V}}=\sum\limits_{\lambda\in\Omega_T}P_{\lambda}$.
\end{enumerate}
For $\mathcal{V}=\{\mathbf{0}\}$, the spectral set and projection family are empty, and both sums are the unique zero operator.
\end{definition}

\begin{theorem} \label{thm:finite-spectral-functional-calculus-normal} \hlt{Finite Spectral Functional Calculus for Normal Operators}\\
Let $\mathcal{V}$ be a finite-dimensional complex inner-product space, and let $T\in\LL(\mathcal{V})$ be normal.\\
Define $\Phi_T:\mathbb{C}^{\Omega_T}\to\LL(\mathcal{V})$ by $\Phi_T(f)\defeq f(T)\defeq\sum\limits_{\lambda\in\Omega_T}f(\lambda)P_{\lambda}$.
\begin{enumerate}[label=(\roman*)]
\setlength{\itemsep}{0pt}
\item $(\forall f\in\mathbb{C}^{\Omega_T})(\forall p\in\mathbb{C}[z])\,[p|_{\Omega_T}=f\Longrightarrow p(T)=f(T)]$.
\item $(\forall f\in\mathbb{C}^{\Omega_T})(\exists p_f\in\mathbb{C}[z])\,[p_f|_{\Omega_T}=f\land p_f(T)=f(T)]$.
\end{enumerate}
\begin{proof}
If $\Omega_T=\emptyset$, then $\mathcal{V}=\{\mathbf{0}\}$ by Theorem \ref{thm:complex-spectral-theorem}(i)--(ii).\\
The zero polynomial represents the unique $f\in\mathbb{C}^{\Omega_T}$, and all operator sums vanish; hence (i)--(ii).\\
Otherwise $\Omega_T\neq\emptyset$, and the following arguments apply.
\begin{enumerate}[label=(\roman*)]
\setlength{\itemsep}{0pt}
\item Fix $f\in\mathbb{C}^{\Omega_T}$ and $p\in\mathbb{C}[z]$ with $p|_{\Omega_T}=f$.\\
Then $p(T)=\sum\limits_{\lambda\in\Omega_T}p(\lambda)P_{\lambda}=\sum\limits_{\lambda\in\Omega_T}f(\lambda)P_{\lambda}=f(T)$ by Corollary \ref{cor:normal-spectral-resolution-functional-calculus}(i) and the restriction identity.
\item Fix $f\in\mathbb{C}^{\Omega_T}$. For $\lambda\in\Omega_T$, set $\ell_{\lambda}(z)\defeq\prod\limits_{\mu\in\Omega_T\setminus\{\lambda\}}\frac{z-\mu}{\lambda-\mu}$.\\
Set $p_f(z)\defeq\sum\limits_{\lambda\in\Omega_T}f(\lambda)\ell_{\lambda}(z)$. For $\lambda_0\in\Omega_T$, $(\forall\lambda\in\Omega_T)\,\ell_{\lambda}(\lambda_0)=\delta_{\lambda,\lambda_0}$ by distinctness.\\
Thus $p_f(\lambda_0)=f(\lambda_0)$, so $p_f|_{\Omega_T}=f$; moreover, $p_f(T)=f(T)$ by (i).
\end{enumerate}
\hlt{Structural Commentary}: Polynomial spectral evaluation feeds an explicit Lagrange interpolation on $\Omega_T$.\\
\hlt{Key Insight}: $\ell_{\lambda}(\mu)=\delta_{\lambda,\mu}$ encodes arbitrary finite spectral data.
\end{proof}
\end{theorem}

\begin{theorem} \label{thm:normal-spectral-functional-calculus-rules} \hlt{Rules for Normal Spectral Functions}\\
Let $\mathcal{V}$ be a finite-dimensional complex inner-product space, and let $T\in\LL(\mathcal{V})$ be normal.\\
For $f,g:\Omega_T\to\mathbb{C}$, the following identities and properties hold:
\begin{enumerate}[label=(\roman*)]
\setlength{\itemsep}{0pt}
\item $(f+g)(T)=f(T)+g(T)$;
\item $(fg)(T)=f(T)g(T)$;
\item $f(T)^*=\overline{f}(T)$, where $(\forall\lambda\in\Omega_T)\,\overline{f}(\lambda)\defeq\overline{f(\lambda)}$;
\item $f(T)$ is normal;
\item $\operatorname{spec}(f(T))=f(\Omega_T)$.
\end{enumerate}
\begin{proof}
If $\mathcal{V}=\{\mathbf{0}\}$, all spectral sums and spectra are empty, so (i)--(v) hold. Assume $\mathcal{V}\neq\{\mathbf{0}\}$.
\begin{enumerate}[label=(\roman*)]
\setlength{\itemsep}{0pt}
\item $(f+g)(T)=\sum\limits_{\lambda\in\Omega_T}(f(\lambda)+g(\lambda))P_{\lambda}=f(T)+g(T)$ by the definition of the spectral functional calculus.
\item The relations $P_{\lambda}P_{\mu}=0$ for $\lambda\neq\mu$ and $P_{\lambda}^2=P_{\lambda}$ hold by Theorem \ref{thm:spectral-resolution-normal-operator}(i)--(ii). Hence $f(T)g(T)=\sum\limits_{\lambda\in\Omega_T}f(\lambda)g(\lambda)P_{\lambda}=(fg)(T)$.
\item $(\forall\lambda\in\Omega_T)\,P_{\lambda}^*=P_{\lambda}$ by Theorem \ref{thm:orthogonal-projection-self-adjoint-idempotent}. Thus $f(T)^*=\sum\limits_{\lambda\in\Omega_T}\overline{f(\lambda)}P_{\lambda}=\overline{f}(T)$ by Proposition \ref{prop:basic-adjoint-identities}(i).
\item $f(T)^*f(T)=(\overline{f}f)(T)=(f\overline{f})(T)=f(T)f(T)^*$ by (ii)--(iii). Therefore $f(T)$ is normal by Definition \ref{def:normal-operator}(i).
\item Set $n\defeq\dim(\mathcal{V})$. Fix an orthonormal eigenbasis $(\mathbf{e}_j)_{j=1}^{n}$ for $T$ by Theorem \ref{thm:complex-spectral-theorem}(i)--(ii), and write $T\mathbf{e}_j=\lambda_j\mathbf{e}_j$.\\
Then $f(T)\mathbf{e}_j=f(\lambda_j)\mathbf{e}_j$ by the spectral resolution. Thus $f(T)$ is diagonal, with diagonal-value set $f(\Omega_T)$, so $\operatorname{spec}(f(T))=f(\Omega_T)$.
\end{enumerate}
\hlt{Structural Commentary}: Projection arithmetic transfers pointwise algebra and conjugation before the eigenbasis reads the spectrum.\\
\hlt{Key Insight}: Every $E_{\lambda}(T)$ converts $f(T)$ into multiplication by $f(\lambda)$.
\end{proof}
\end{theorem}

\begin{definition} \label{def:positive-operator} \hlt{Positive Operator}\\
Let $\mathcal{V}$ be finite-dimensional inner-product space over $\mathbb{F}\in\{\mathbb{R},\mathbb{C}\}$. For $T\in\LL(\mathcal{V})$, $T\geq0$, $T$ \hlt{positive} if $T=T^*$ and $(\forall\mathbf{v}\in\mathcal{V})\,\innerprod{T\mathbf{v},\mathbf{v}}\geq0$.
\end{definition}

\begin{definition} \label{def:square-root-operator} \hlt{Square Root of an Operator}\\
Let $\mathcal{V}$ be a finite-dimensional inner-product space over $\mathbb{F}\in\{\mathbb{R},\mathbb{C}\}$, $T\in\LL(\mathcal{V})$. An operator $R\in\LL(\mathcal{V})$ is a \hlt{square root} of $T$ if $R^2=T$.
\end{definition}

\begin{theorem} \label{thm:positive-operator-characterization} \hlt{Characterization of Positive Operators}\\
Let $\mathcal{V}$ be a finite-dimensional inner-product space over $\mathbb{F}\in\{\mathbb{R},\mathbb{C}\}$, $T\in\LL(\mathcal{V})$. The following are equivalent:
\begin{enumerate}[label=(\roman*)]
\setlength{\itemsep}{0pt}
\item $T$ is positive;
\item $T=T^*$ and $\operatorname{spec}(T)\subseteq[0,\infty)$;
\item $(\exists S\in\LL(\mathcal{V}))\,[S\geq0\land S^2=T]$;
\item $(\exists R\in\LL(\mathcal{V}))\,[R=R^*\land R^2=T]$;
\item $(\exists R\in\LL(\mathcal{V}))\,T=R^*R$.
\end{enumerate}
\begin{proof}
If $\mathcal{V}=\{\mathbf{0}\}$, its unique operator satisfies all five conditions. Assume $\mathcal{V}\neq\{\mathbf{0}\}$.
\begin{enumerate}[label=(\roman*)]
\setlength{\itemsep}{0pt}
\item[$(i)$] $\Rightarrow (ii)$ The identity $T=T^*$ holds by Definition \ref{def:positive-operator}. Fix $\lambda\in\operatorname{spec}(T)$ and $\mathbf{v}\in E_{\lambda}(T)\setminus\{\mathbf{0}\}$.\\
Then $\lambda\in\mathbb{R}$ by Theorem \ref{thm:self-adjoint-eigenvalues-real}, while $0\leq\innerprod{T\mathbf{v},\mathbf{v}}=\lambda\norm{\mathbf{v}}^2$. Hence $\lambda\geq0$.
\item[$(ii)$] $\Rightarrow (iii)$ Set $n\defeq\dim(\mathcal{V})$. In the complex case, $T^*T=T^2=TT^*$ by $T=T^*$, so $T$ is normal.\\
Fix an orthonormal eigenbasis $(\mathbf{e}_j)_{j=1}^{n}$ by Theorem \ref{thm:real-spectral-theorem}(i)--(ii) when $\mathbb{F}=\mathbb{R}$ and by Theorem \ref{thm:complex-spectral-theorem}(i)--(ii) when $\mathbb{F}=\mathbb{C}$.\\
Write $T\mathbf{e}_j=\lambda_j\mathbf{e}_j$, where $(\forall j\in\{1,\ldots,n\})\,\lambda_j\geq0$. Fix $S\in\LL(\mathcal{V})$ satisfying $(\forall j\in\{1,\ldots,n\})\,S\mathbf{e}_j=\sqrt{\lambda_j}\mathbf{e}_j$ by Proposition \ref{prop:linear-map-determined-by-basis}.\\
For $\mathbf{u}=\sum\limits_{j=1}^{n}c_j\mathbf{e}_j$ and $\mathbf{v}=\sum\limits_{j=1}^{n}d_j\mathbf{e}_j$, $\innerprod{S\mathbf{u},\mathbf{v}}=\sum\limits_{j=1}^{n}\sqrt{\lambda_j}c_j\overline{d_j}=\innerprod{\mathbf{u},S\mathbf{v}}$.\\
$S=S^*$ by the preceding adjoint identity, and $\innerprod{S\mathbf{v},\mathbf{v}}=\sum\limits_{j=1}^{n}\sqrt{\lambda_j}\abs{d_j}^2\geq0$. Hence $S\geq0$ by Definition \ref{def:positive-operator}.\\
The identity $S^2=T$ holds by Proposition \ref{prop:linear-map-determined-by-basis} and the equality $(\forall j\in\{1,\ldots,n\})\,S^2\mathbf{e}_j=T\mathbf{e}_j$.
\item[$(iii)$] $\Rightarrow (iv)$ The relation $S=S^*$ holds by $S\geq0$ and Definition \ref{def:positive-operator}, so $S$ is a self-adjoint square root.
\item[$(iv)$] $\Rightarrow (v)$ The identity $T=R^*R$ holds by $T=R^2$ and $R=R^*$.
\item[$(v)$] $\Rightarrow (i)$ The identity $T^*=(R^*R)^*=R^*(R^*)^*=R^*R=T$ holds by Proposition \ref{prop:basic-adjoint-identities}(ii)--(iii).\\
The inequality $(\forall\mathbf{v}\in\mathcal{V})\,\innerprod{T\mathbf{v},\mathbf{v}}=\norm{R\mathbf{v}}^2\geq0$ holds by $T=R^*R$ and the defining adjoint identity.\\
Thus $T\geq0$ by Definition \ref{def:positive-operator} and $T^*=T$.
\end{enumerate}
\hlt{Structural Commentary}: Spectral coordinates construct the positive root, and adjoint factorization closes the implication cycle.\\
\hlt{Key Insight}: $\lambda_j\geq0$ permits the canonical multiplier $\sqrt{\lambda_j}$.
\end{proof}
\end{theorem}

\begin{proposition} \label{prop:positive-square-root-unique} \hlt{Unique Positive Square Root}\\
Let $\mathcal{V}$ be a finite-dimensional inner-product space over $\mathbb{F}\in\{\mathbb{R},\mathbb{C}\}$. For $T\in\LL(\mathcal{V})$ with $T\geq0$, $(\exists!R\in\LL(\mathcal{V}))\,[R\geq0\land R^2=T]$.
\begin{proof}
\begin{enumerate}[label=(\roman*)]
\setlength{\itemsep}{0pt}
\item \hlt{Existence and Degenerate Case}: A positive square root exists by Theorem \ref{thm:positive-operator-characterization}(i),(iii).\\
If $\mathcal{V}=\{\mathbf{0}\}$, its unique operator is its unique positive square root. Assume $\mathcal{V}\neq\{\mathbf{0}\}$.
\item \hlt{Spectral Reduction}: Let $\lambda_1,\ldots,\lambda_m$ be the distinct eigenvalues of $T$.\\
Then $T=T^*$ and $(\forall j\in\{1,\ldots,m\})\,\lambda_j\geq0$ by Theorem \ref{thm:positive-operator-characterization}(i),(ii). In the complex case, $T^*T=T^2=TT^*$, so $T$ is normal.\\
Hence $\mathcal{V}=E_{\lambda_1}(T)\odot\cdots\odot E_{\lambda_m}(T)$ by Theorem \ref{thm:real-spectral-theorem}(i)--(ii) when $\mathbb{F}=\mathbb{R}$ and by Theorem \ref{thm:complex-spectral-theorem}(i)--(ii) when $\mathbb{F}=\mathbb{C}$.
\item \hlt{Restriction to One Eigenspace}: Let $R$ be a positive square root of $T$.\\
The identity $RT=R^3=TR$ holds by $R^2=T$, so $(\forall j\in\{1,\ldots,m\})\,E_{\lambda_j}(T)\in\operatorname{Inv}(R)$.\\
For $j\in\{1,\ldots,m\}$, set $R_j\defeq R|_{E_{\lambda_j}(T)}$. Then $R_j=R_j^*$ by Proposition \ref{prop:self-adjoint-invariant-subspaces}(ii), and $(\forall\mathbf{x}\in E_{\lambda_j}(T))\,\innerprod{R_j\mathbf{x},\mathbf{x}}\geq0$ by $R\geq0$.\\
Thus $R_j\geq0$ by Definition \ref{def:positive-operator}, while $R_j^2=\lambda_jI_{E_{\lambda_j}(T)}$ by $R^2=T$. In the complex case, $R_j^*R_j=R_j^2=R_jR_j^*$.\\
Fix an orthonormal eigenbasis of $R_j$ by Theorem \ref{thm:real-spectral-theorem}(i)--(ii) when $\mathbb{F}=\mathbb{R}$ and by Theorem \ref{thm:complex-spectral-theorem}(i)--(ii) when $\mathbb{F}=\mathbb{C}$.\\
Every corresponding eigenvalue $\rho$ satisfies $\rho\geq0$ by Theorem \ref{thm:positive-operator-characterization}(i),(ii) and $\rho^2=\lambda_j$ by the restriction identity.\\
Therefore $\rho=\sqrt{\lambda_j}$, so $R_j=\sqrt{\lambda_j}I_{E_{\lambda_j}(T)}$ by Proposition \ref{prop:linear-map-determined-by-basis}.
\item \hlt{Uniqueness}: Every positive square root $R$ satisfies $(\forall j\in\{1,\ldots,m\})\,R|_{E_{\lambda_j}(T)}=\sqrt{\lambda_j}I_{E_{\lambda_j}(T)}$ by (iii).\\
Any two positive square roots agree on every summand by this identity and are therefore equal on $\mathcal{V}$ by the decomposition in (ii).
\end{enumerate}
\hlt{Structural Commentary}: Spectral reduction localizes uniqueness to the restriction equation $R_j^2=\lambda_jI$.\\
\hlt{Key Insight}: Positivity selects $\sqrt{\lambda_j}$ rather than $-\sqrt{\lambda_j}$.
\end{proof}
\end{proposition}

\begin{definition} \label{def:positive-square-root-notation} \hlt{Positive Square Root Notation}\\
Let $\mathcal{V}$ be finite-dimensional inner-product space over $\mathbb{F}\in\{\mathbb{R},\mathbb{C}\}$. For $T\in\LL(\mathcal{V})$, $T\geq0$, $\sqrt{T}$ denotes its unique positive square root.
\end{definition}

\begin{corollary} \label{cor:spectral-functional-calculus-basic-functions} \hlt{Basic Functions from the Spectral Calculus}\\
Let $\mathcal{V}$ be a finite-dimensional complex inner-product space, and let $T\in\LL(\mathcal{V})$ be normal.
\begin{enumerate}[label=(\roman*)]
\setlength{\itemsep}{0pt}
\item $T^*=\sum\limits_{\lambda\in\Omega_T}\overline{\lambda}P_{\lambda}$.
\item If $0\notin\Omega_T$, then $T^{-1}=\sum\limits_{\lambda\in\Omega_T}\lambda^{-1}P_{\lambda}$.
\item If $T\geq0$, then $\sqrt{T}=\sum\limits_{\lambda\in\Omega_T}\sqrt{\lambda}P_{\lambda}$.
\end{enumerate}
\begin{proof}
\begin{enumerate}[label=(\roman*)]
\setlength{\itemsep}{0pt}
\item $T^*=\sum\limits_{\lambda\in\Omega_T}\overline{\lambda}P_{\lambda}$ by Theorem \ref{thm:normal-spectral-functional-calculus-rules}(iii) with the identity function on $\Omega_T$.
\item Set $f(\lambda)\defeq\lambda^{-1}$. Then $Tf(T)=f(T)T=\sum\limits_{\lambda\in\Omega_T}P_{\lambda}=I_{\mathcal{V}}$ by the projection product relations and the resolution of $I_{\mathcal{V}}$.\\
Thus $T^{-1}=f(T)=\sum\limits_{\lambda\in\Omega_T}\lambda^{-1}P_{\lambda}$ by uniqueness of the inverse.
\item The inclusion $\Omega_T\subseteq[0,\infty)$ holds by Theorem \ref{thm:positive-operator-characterization}(i),(ii). Set $g(\lambda)\defeq\sqrt{\lambda}$.\\
Then $g(T)^*=g(T)$, $g(T)^2=T$, and $\operatorname{spec}(g(T))=g(\Omega_T)\subseteq[0,\infty)$ by Theorem \ref{thm:normal-spectral-functional-calculus-rules}(ii)--(v).\\
Hence $g(T)\geq0$ by Theorem \ref{thm:positive-operator-characterization}(ii),(i), and $\sqrt{T}=g(T)=\sum\limits_{\lambda\in\Omega_T}\sqrt{\lambda}P_{\lambda}$ by Proposition \ref{prop:positive-square-root-unique}.
\end{enumerate}
\hlt{Structural Commentary}: Conjugation, inversion, and positive square roots transfer pointwise through the spectral calculus.\\
\hlt{Key Insight}: Each scalar function acts independently on every spectral block.
\end{proof}
\end{corollary}

\begin{theorem} \label{thm:normal-spectral-projection-commutativity} \hlt{Commutativity via Spectral Projections}\\
Let $\mathcal{V}$ be a finite-dimensional complex inner-product space, and let $T,S\in\LL(\mathcal{V})$ be normal. Write $T=\sum\limits_{\lambda\in\Omega_T}\lambda P_{\lambda}$ and $S=\sum\limits_{\mu\in\Omega_S}\mu Q_{\mu}$.
\begin{enumerate}[label=(\roman*)]
\setlength{\itemsep}{0pt}
\item For $R\in\LL(\mathcal{V})$, $RT=TR\Longleftrightarrow(\forall\lambda\in\Omega_T)\,RP_{\lambda}=P_{\lambda}R$.
\item $TS=ST\Longleftrightarrow(\forall\lambda\in\Omega_T)(\forall\mu\in\Omega_S)\,P_{\lambda}Q_{\mu}=Q_{\mu}P_{\lambda}$.
\item If $f:\Omega_T\to\mathbb{C}$ and $g:\Omega_S\to\mathbb{C}$ are injective, then $f(T)g(S)=g(S)f(T)\Longleftrightarrow TS=ST$.
\end{enumerate}
\begin{proof}
If $\mathcal{V}=\{\mathbf{0}\}$, both sides of every equivalence hold. Assume $\mathcal{V}\neq\{\mathbf{0}\}$.
\begin{enumerate}[label=(\roman*)]
\setlength{\itemsep}{0pt}
\item[$(i)$] ($\Rightarrow$) If $RT=TR$, then $(\forall p\in\mathbb{C}[z])\,Rp(T)=p(T)R$ by induction on powers and linearity.\\
Each $P_{\lambda}$ is a polynomial in $T$ by Corollary \ref{cor:normal-spectral-resolution-functional-calculus}(ii), so $(\forall\lambda\in\Omega_T)\,RP_{\lambda}=P_{\lambda}R$.
\item[$(i)$] ($\Leftarrow$) If $(\forall\lambda\in\Omega_T)\,RP_{\lambda}=P_{\lambda}R$, then $RT=\sum\limits_{\lambda\in\Omega_T}\lambda RP_{\lambda}=\sum\limits_{\lambda\in\Omega_T}\lambda P_{\lambda}R=TR$.
\item[$(ii)$] ($\Rightarrow$) If $TS=ST$, then $(\forall\lambda\in\Omega_T)\,SP_{\lambda}=P_{\lambda}S$ by (i) with $R=S$. Hence $(\forall\lambda\in\Omega_T)(\forall\mu\in\Omega_S)\,P_{\lambda}Q_{\mu}=Q_{\mu}P_{\lambda}$ by (i) for $S$ with $R=P_{\lambda}$.
\item[$(ii)$] ($\Leftarrow$) If all spectral projections commute pairwise, then $TS=\sum\limits_{\lambda\in\Omega_T,\,\mu\in\Omega_S}\lambda\mu P_{\lambda}Q_{\mu}=\sum\limits_{\lambda\in\Omega_T,\,\mu\in\Omega_S}\lambda\mu Q_{\mu}P_{\lambda}=ST$.
\item[$(iii)$] ($\Rightarrow$) Suppose $f(T)g(S)=g(S)f(T)$. Then $f(T)$ and $g(S)$ are normal by Theorem \ref{thm:normal-spectral-functional-calculus-rules}(iv).\\
Spectral coefficients are distinct by injectivity, so spectral projections are $(P_{\lambda})_{\lambda\in\Omega_T}$ and $(Q_{\mu})_{\mu\in\Omega_S}$ by Theorem \ref{thm:spectral-resolution-normal-operator}(ii).\\
These projections commute by (ii) for $f(T)$ and $g(S)$; hence $TS=ST$ by (ii) for $T$ and $S$.
\item[$(iii)$] ($\Leftarrow$) If $TS=ST$, then all $P_{\lambda}$ commute with all $Q_{\mu}$ by (ii).\\
Thus $f(T)g(S)=\sum\limits_{\lambda\in\Omega_T,\,\mu\in\Omega_S}f(\lambda)g(\mu)P_{\lambda}Q_{\mu}=\sum\limits_{\lambda\in\Omega_T,\,\mu\in\Omega_S}f(\lambda)g(\mu)Q_{\mu}P_{\lambda}=g(S)f(T)$.
\end{enumerate}
\hlt{Structural Commentary}: Polynomial recovery transfers operator commutation to pairwise commutation of spectral projections.\\
\hlt{Key Insight}: Commutation is equivalent to preserving every spectral block.
\end{proof}
\end{theorem}

\begin{corollary} \label{cor:commuting-normal-pair-generated-by-one-normal-operator} \hlt{A Commuting Normal Pair Comes from One Normal Operator}\\
Let $\mathcal{V}$ be a finite-dimensional complex inner-product space, and let $T,S\in\LL(\mathcal{V})$ be normal.\\
Let $\mathbb{C}[T,S]$ denote the finite sums of $c_{a,b}T^aS^b$ with $c_{a,b}\in\mathbb{C}$ and $a,b\in\mathbb{Z}_{\geq0}$.\\
Then $TS=ST\Longleftrightarrow(\exists H\in\mathbb{C}[T,S])(\exists p,q\in\mathbb{C}[z])\,[H^*H=HH^*\land T=p(H)\land S=q(H)]$.
\begin{proof}
\begin{enumerate}[label=\roman*.]
\setlength{\itemsep}{0pt}
\item[$(\Rightarrow)$] Suppose $TS=ST$. If $\mathcal{V}=\{\mathbf{0}\}$, the choices $H=T=S=0$ and $p=q=0$ satisfy the conclusion. Assume $\mathcal{V}\neq\{\mathbf{0}\}$.
\begin{enumerate}[label=(\alph*)]
\setlength{\itemsep}{0pt}
\item \hlt{Joint Spectral Blocks}: Write $T=\sum\limits_{\lambda\in\Omega_T}\lambda P_{\lambda}$ and $S=\sum\limits_{\mu\in\Omega_S}\mu Q_{\mu}$.\\
The projections commute by Theorem \ref{thm:normal-spectral-projection-commutativity}(ii). Set $\mathcal{I}\defeq\{(\lambda,\mu)\in\Omega_T\times\Omega_S:P_{\lambda}Q_{\mu}\neq0\}$.\\
For $(\lambda,\mu)\in\mathcal{I}$, $(P_{\lambda}Q_{\mu})^2=P_{\lambda}Q_{\mu}$ and $(P_{\lambda}Q_{\mu})^*=P_{\lambda}Q_{\mu}$ by projection commutation and self-adjoint idempotence.\\
Thus $P_{\lambda}Q_{\mu}$ is an orthogonal projection by Theorem \ref{thm:orthogonal-projection-self-adjoint-idempotent}.\\
Distinct joint projections are orthogonal by pairwise annihilation and commutation.\\
After insertion of the zero products, $\sum\limits_{(\lambda,\mu)\in\mathcal{I}}P_{\lambda}Q_{\mu}=\sum\limits_{\lambda\in\Omega_T,\,\mu\in\Omega_S}P_{\lambda}Q_{\mu}$.\\
The latter sum equals $(\sum\limits_{\lambda\in\Omega_T}P_{\lambda})(\sum\limits_{\mu\in\Omega_S}Q_{\mu})=I_{\mathcal{V}}$ by the spectral resolutions.
\item \hlt{Separating Operator}: Choose pairwise distinct $\alpha_{\lambda,\mu}\in\mathbb{C}$ for $(\lambda,\mu)\in\mathcal{I}$ and set $H\defeq\sum\limits_{(\lambda,\mu)\in\mathcal{I}}\alpha_{\lambda,\mu}P_{\lambda}Q_{\mu}$.\\
Then $H$ is normal, with $\operatorname{spec}(H)=\{\alpha_{\lambda,\mu}:(\lambda,\mu)\in\mathcal{I}\}$, by Theorem \ref{thm:spectral-resolution-normal-operator}(i)--(ii).
\item \hlt{Generated Algebra}: Each $P_{\lambda}$ is a polynomial in $T$ and each $Q_{\mu}$ is a polynomial in $S$ by Corollary \ref{cor:normal-spectral-resolution-functional-calculus}(ii).\\
Therefore $H$ is a finite sum of terms $c_{a,b}T^aS^b$, so $H\in\mathbb{C}[T,S]$.
\item \hlt{Recovery of the Pair}: Define $f,g:\operatorname{spec}(H)\to\mathbb{C}$ by $f(\alpha_{\lambda,\mu})\defeq\lambda$ and $g(\alpha_{\lambda,\mu})\defeq\mu$.\\
Fix $p,q\in\mathbb{C}[z]$ satisfying $p|_{\operatorname{spec}(H)}=f$ and $q|_{\operatorname{spec}(H)}=g$ by Theorem \ref{thm:finite-spectral-functional-calculus-normal}(ii).\\
Then $p(H)=\sum\limits_{(\lambda,\mu)\in\mathcal{I}}\lambda P_{\lambda}Q_{\mu}=T$ by the spectral resolution of $T$.\\
The identity $q(H)=\sum\limits_{(\lambda,\mu)\in\mathcal{I}}\mu P_{\lambda}Q_{\mu}=S$ holds by the spectral resolution of $S$.
\end{enumerate}
\item[$(\Leftarrow)$] Suppose $T=p(H)$ and $S=q(H)$. Then $p(H)q(H)=q(H)p(H)$ by $H^aH^b=H^{a+b}=H^bH^a$ and distributivity, so $TS=ST$.
\end{enumerate}
\hlt{Structural Commentary}: Joint spectral blocks build a separating normal generator before polynomial recovery of $T$ and $S$.\\
\hlt{Key Insight}: Distinct labels encode both eigenvalue coordinates on each joint block.
\end{proof}
\end{corollary}
